\documentclass[aspectratio=169,10pt,spanish]{beamer}

\input{../../../_preambulo_beamer.tex}
\input{../../../_comandos_beamer.tex}

\selectlanguage{spanish}

\title{MA1001B - Modelación Estadística para la Toma de Decisiones}
\subtitle{Lección 04: Detección de anomalías y boxplots}
\author{}
\institute{Tecnol\'ogico de Monterrey}
\date{}

\begin{document}

\begin{frame}[plain]
  \vspace{1.2cm}
  {\Large\bfseries MA1001B - Modelación Estadística para la Toma de Decisiones\par}
  \vspace{0.7cm}
  {\large Lección 04: Detección de anomalías y boxplots\par}
  \vspace{0.7cm}
  {\color{TecRojo}\rule{\textwidth}{1.2pt}}
  \vspace{0.8cm}

  Facultad MA1001B\\[0.3cm]
  Periodo\\[0.3cm]
  Tecnol\'ogico de Monterrey
\end{frame}

\begin{frame}{La pregunta de cribado}
  \begin{alertblock}{Pregunta guía}
    ¿Qué tiempos de procesamiento merecen investigación y qué puede decirnos
    realmente una bandera estadística?
  \end{alertblock}

  \vspace{0.5em}
  Al final de esta lección, usted puede:
  \begin{itemize}
    \item enunciar cuándo es apropiada la regla empírica;
    \item aplicar la desigualdad de Chebyshov sin supuesto de forma;
    \item calcular e interpretar las cercas 1.5-IQR de un boxplot;
    \item separar una bandera de cribado de una conclusión sobre su causa.
  \end{itemize}

  \vfill
  \scriptsize Los 600 registros de pedidos son sintéticos. No se usan datos reales de empresa.
\end{frame}

\begin{frame}{Tres reglas de distancia}
  \small
  \begin{center}
    \begin{tabular}{p{0.22\textwidth}p{0.33\textwidth}p{0.34\textwidth}}
      \toprule
      \textbf{Regla} & \textbf{Condición} & \textbf{Qué aporta} \\
      \midrule
      Regla empírica & Datos simétricos con forma de campana & Cobertura aproximada dentro de $1s$, $2s$ y $3s$ \\
      Chebyshov & Cualquier distribución, con $k>1$ & Cota inferior $1-1/k^2$ dentro de $ks$ \\
      Cercas de boxplot & Datos cuantitativos ordenados & Valores más allá de las cercas 1.5-IQR \\
      \bottomrule
    \end{tabular}
  \end{center}

  \vspace{0.7em}
  \begin{exampleblock}{Límite compartido}
    Cada regla describe posición estadística. La calidad del registro, el
    contexto operativo y la importancia de negocio requieren evidencia
    adicional.
  \end{exampleblock}
\end{frame}

\begin{frame}{Paso 1: Regla empírica bajo su condición}
  \begin{columns}[T]
    \begin{column}{0.37\textwidth}
      \small
      Tiempos sintéticos con forma de campana:

      \vspace{0.4em}
      $n=600$ pedidos\\
      Media $=44.87$ min\\
      $s$ muestral $=4.87$ min

      \vspace{0.7em}
      \begin{tabular}{lr}
        \toprule
        \textbf{Intervalo} & \textbf{Observado} \\
        \midrule
        Dentro de $1s$ & 69.17\% \\
        Dentro de $2s$ & 94.83\% \\
        Dentro de $3s$ & 99.50\% \\
        \bottomrule
      \end{tabular}

      \vspace{0.5em}
      Estos valores respaldan el patrón aproximado de 68\%, 95\% y más de
      99\% para esta muestra.
    \end{column}
    \begin{column}{0.63\textwidth}
      \centering
      \includegraphics[width=\textwidth]{../figuras/deteccion_anomalias_01_regla_empirica.png}
    \end{column}
  \end{columns}
\end{frame}

\begin{frame}{Paso 2: Una cota inferior sin forma}
  \begin{columns}[T]
    \begin{column}{0.39\textwidth}
      \small
      Los retrasos de cola crean una cola derecha:

      \vspace{0.3em}
      59 pedidos retrasados\\
      Media $=45.79$ min\\
      $s$ muestral $=5.84$ min\\
      Asimetría $=0.74$

      \vspace{0.5em}
      \[
        \text{Mínimo dentro de }ks=1-\frac{1}{k^2}
      \]

      \begin{tabular}{lrr}
        \toprule
        & \textbf{Cota} & \textbf{Observado} \\
        \midrule
        Dentro de $2s$ & 75.00\% & 95.33\% \\
        Dentro de $3s$ & 88.89\% & 98.67\% \\
        \bottomrule
      \end{tabular}
    \end{column}
    \begin{column}{0.61\textwidth}
      \centering
      \includegraphics[width=0.97\textwidth]{../figuras/deteccion_anomalias_02_chebyshev.png}
    \end{column}
  \end{columns}
\end{frame}

\begin{frame}{Paso 3: Candidatos a investigación del boxplot}
  \small
  \begin{columns}[T]
    \begin{column}{0.36\textwidth}
      \[
        \mathrm{IQR}=Q_3-Q_1=6.71
      \]
      $Q_1=42.19$ min\\
      $Q_3=48.89$ min\\
      Cerca inferior $=32.13$ min\\
      Cerca superior $=58.96$ min

      \vspace{0.8em}
      \begin{alertblock}{Resultado de cribado verificado}
        La regla marca 24 pedidos. Captura los cuatro extremos inyectados y
        también marca 20 observaciones más.
      \end{alertblock}
    \end{column}
    \begin{column}{0.64\textwidth}
      \centering
      \includegraphics[width=\textwidth]{../figuras/deteccion_anomalias_03_banderas_boxplot.png}
    \end{column}
  \end{columns}
\end{frame}

\begin{frame}{De la bandera a la decisión}
  \small
  \begin{center}
    \begin{tabular}{p{0.21\textwidth}p{0.69\textwidth}}
      \toprule
      \textbf{Etapa} & \textbf{Evidencia necesaria} \\
      \midrule
      Cribado estadístico & Una regla enunciada, un cálculo reproducible y el valor original conservado \\
      Revisión del registro & Unidad, marca de tiempo, identificador e historial de transformaciones correctos \\
      Contexto del proceso & Eventos relacionados, condiciones operativas y registros comparables \\
      Resolución & Razón documentada para retener, corregir, excluir o modelar por separado \\
      \bottomrule
    \end{tabular}
  \end{center}

  \vspace{0.8em}
  \begin{exampleblock}{Qué significan las otras 20 banderas}
    Cumplen la misma regla 1.5-IQR y carecen de la etiqueta de construcción
    sintética conocida. Sus causas permanecen desconocidas. Llamarlas
    errores o falsos positivos exigiría evidencia que el cribado
    estadístico no contiene.
  \end{exampleblock}
\end{frame}

\begin{frame}{Verificación de conceptos}
  \begin{enumerate}
    \item ¿Por qué la regla empírica sería cuestionable para la muestra con
      cola derecha del Paso 2?
    \item Calcule la cota inferior de Chebyshov para $k=2$ e interprétela.
    \item Con $Q_1=42.19$ y $Q_3=48.89$, reproduzca la cerca superior.
    \item ¿Qué evidencia buscaría antes de modificar un registro marcado?
  \end{enumerate}

  \vspace{0.5em}
  \begin{block}{Conexión con la Lección 05}
    La notación de conjuntos y los espacios muestrales ofrecen un lenguaje
    preciso para eventos. Ayudarán a describir qué registros satisfacen una
    regla y cómo se combinan los eventos.
  \end{block}

  \vfill
  \scriptsize Fuente: OpenStax, \textit{Introductory Business Statistics 2e},
  Capítulo 2, Secciones 2.1, 2.2 y 2.7. CC BY-NC-SA 4.0.
\end{frame}

\appendix



\end{document}
